% !TeX program = xelatex
\documentclass[12pt,a4paper]{article}

\usepackage[margin=1in]{geometry}
\usepackage{setspace}
\usepackage{booktabs}
\usepackage{amsmath,amssymb}
\usepackage{hyperref}
\usepackage{graphicx}
\usepackage{float}
\usepackage{pgfplots}
\pgfplotsset{compat=1.18}
\setlength{\emergencystretch}{3em}

\title{Introducing Relay-Expert Distillation: Training Update and Early Results}
\author{}
\date{\today}

\begin{document}

\maketitle

\begin{abstract}
We are sharing an early training update for Relay-Expert Distillation, a routing-based specialist framework designed to improve capability transfer while preserving optimization stability. Starting from the raw log file \url{figures/training_metrics.csv}, we cleaned malformed records and rebuilt a reliable metrics stream in \url{figures/training_metrics_clean.csv}. Across steps 5710 to 13400, loss decreases from 6.2683 to 4.5531, perplexity decreases from 527.60 to 94.93, and the best checkpoint appears at step 13160 (loss=3.8976, ppl=49.29). The full run remains stable and still improves near the end, suggesting additional headroom for short continuation training.
\end{abstract}

\section{Overview}
This release summarizes what happened during training, what signals matter most, and what we plan to do next. The outcome is encouraging: convergence is consistent, optimization is well-behaved, and late-stage metrics indicate the run is not yet fully saturated.

\section{Key Takeaways}
\begin{itemize}
  \item \textbf{Convergence is strong:} both loss and perplexity improve substantially across the full run.
  \item \textbf{Optimization remains stable:} learning-rate decay aligns with metric trends and shows no obvious divergence period.
  \item \textbf{Best quality appears late:} the best loss occurs at step 13160, close to the final step 13400.
  \item \textbf{Efficiency changes are predictable:} throughput gradually decreases from 261 to 235 tok/s.
\end{itemize}

\section{Data and Method}
\subsection{Data Source}
We use the fields \texttt{step, loss, ppl, lr, tok\_s}. The cleaned dataset contains 770 valid records, spanning step range [5710, 13400] with step interval 10.

\subsection{Cleaning Strategy}
The raw CSV contains repeated headers and one malformed concatenated row near the end. We apply three deterministic rules:
\begin{itemize}
  \item remove repeated header rows;
  \item keep only rows matching a five-column numeric schema;
  \item truncate malformed concatenated rows before final validation.
\end{itemize}

\subsection{Metrics}
We track three dimensions:
\begin{itemize}
  \item \textbf{Convergence:} loss and perplexity;
  \item \textbf{Optimization state:} learning rate;
  \item \textbf{System efficiency:} token throughput (tok/s).
\end{itemize}

\section{Training Results}
\subsection{Summary Statistics}
\begin{table}[H]
\centering
\caption{Core metrics from step 5710 to step 13400}
\begin{tabular}{lcccc}
\toprule
Metric & Start (step 5710) & End (step 13400) & Best Value & Best Step \\
\midrule
Loss & 6.2683 & 4.5531 & 3.8976 & 13160 \\
PPL & 527.60 & 94.93 & 49.29 & 13160 \\
LR & $2.96\times10^{-4}$ & $2.61\times10^{-4}$ & N/A & N/A \\
tok/s & 261 & 235 & N/A & N/A \\
\bottomrule
\end{tabular}
\end{table}

Over the final 100 steps, average metrics are loss=4.4312, ppl=87.4712, and tok/s=235.18.

\subsection{Loss Trend}
\begin{figure}[H]
\centering
\begin{tikzpicture}
\begin{axis}[
    width=0.9\linewidth,
    height=6cm,
    xlabel={Step},
    ylabel={Loss},
    grid=major,
    line width=0.8pt,
]
\addplot[blue] table[x=step,y=loss,col sep=comma] {figures/training_metrics_clean.csv};
\end{axis}
\end{tikzpicture}
\caption{Loss declines steadily with moderate local variance}
\end{figure}

\subsection{Perplexity Trend}
\begin{figure}[H]
\centering
\begin{tikzpicture}
\begin{axis}[
    width=0.9\linewidth,
    height=6cm,
    xlabel={Step},
    ylabel={Perplexity},
    grid=major,
    line width=0.8pt,
]
\addplot[red!70!black] table[x=step,y=ppl,col sep=comma] {figures/training_metrics_clean.csv};
\end{axis}
\end{tikzpicture}
\caption{Perplexity tracks the same downward trajectory}
\end{figure}

\subsection{Learning Rate and Throughput}
\begin{figure}[H]
\centering
\begin{tikzpicture}
\begin{axis}[
    width=0.9\linewidth,
    height=6cm,
    xlabel={Step},
    ylabel={Learning Rate},
    grid=major,
    line width=0.8pt,
]
\addplot[teal!70!black] table[x=step,y=lr,col sep=comma] {figures/training_metrics_clean.csv};
\end{axis}
\end{tikzpicture}
\caption{Learning-rate decay is smooth across the run}
\end{figure}

\begin{figure}[H]
\centering
\begin{tikzpicture}
\begin{axis}[
    width=0.9\linewidth,
    height=6cm,
    xlabel={Step},
    ylabel={tok/s},
    grid=major,
    line width=0.8pt,
]
\addplot[orange!85!black] table[x=step,y=tok_s,col sep=comma] {figures/training_metrics_clean.csv};
\end{axis}
\end{tikzpicture}
\caption{Throughput gradually decreases as training progresses}
\end{figure}

\section{Interpretation}
Three signals are most important. First, convergence is broad-based: both loss and perplexity improve materially from start to finish. Second, optimization remains stable: we do not observe a sustained divergence regime. Third, the best checkpoint appears near the end of training (step 13160 compared with final step 13400), which indicates additional training headroom.

Based on these signals, the most practical next move is a short continuation phase with stricter validation monitoring. We recommend combining moving-average trend tracking with patience-based early stopping to capture remaining gains while controlling overfitting risk.

\section{What Comes Next}
We plan to extend this release in four directions:
\begin{itemize}
  \item add confidence intervals and smoothing diagnostics to separate trend from noise;
  \item evaluate by task buckets (math, language, logic) to localize specialist bottlenecks;
  \item monitor validation-side indicators jointly with training metrics to detect turning points earlier;
  \item run controlled ablations on routing policy and specialist allocation ratio to isolate architecture effects.
\end{itemize}

The objective is not only better end metrics, but also a clearer causal account of where Relay-Expert Distillation delivers its gains.

\end{document}